\paragraph{The contrast.} A father's childhood layoff moves his sons' and his daughters' pay in opposite directions, as
Figure~\ref{fig:layoffbands} shows. In one reduced specification for both sexes, a father's teen layoff changes his
son's hourly wage between 25--29 and 45--49 by $-0.538^{**}$ and his daughter's by $+0.387^{*}$; within the same person
the slopes are $-0.086$ (s.e.\ 0.096) and $+0.168^{**}$ ($p = 0.025$) (panel B of Table~\ref{tab:layoff_exog}).
Stacking sons and daughters and testing the difference directly, it is 0.925$^{***}$ ($p = 0.001$) in the wage change and
1.294$^{***}$ in the earnings change, and 0.254$^{**}$ ($p = 0.049$) and 0.380$^{**}$ ($p = 0.030$) a decade within
person. The two change differences survive family-wise corrections (adjusted $p$ 0.004 and 0.003 by Holm, 0.018 and
0.011 by Romano--Wolf); the within-person ones are borderline (0.061 and 0.055). With one job-end rule in every wave the change differences are 1.295$^{**}$ and 1.753$^{**}$. The
difference is positive in all three birth cohorts; the daughters' own gain is clear in those born before 1963
(0.221$^{*}$) and from 1976 (0.644$^{***}$) and nil in those born 1963--75.\footnote{Sons and daughters of the same family
share the cluster. 145 daughters with a teen layoff year are observed at 45--49. Without the children of self-employed
fathers the wage-change difference is 0.773$^{**}$ ($p = 0.010$) and the within-person one 0.259$^{**}$ ($p = 0.033$);
holding the father's self-employment fixed, 0.901$^{***}$. Code: \texttt{92.LayoffDaughters.py},
\texttt{107.LayoffSexDifference.py}, \texttt{119.FatherSelfEmployment.py}.} A trait of the father passed on to all his children alike would lower both profiles. Nor is it that a parent's layoff
helps daughters: the layoffs of mothers who head their families lower their daughters' earnings profiles
($-0.561^{**}$ a decade within the daughter; Section~\ref{sec:stability}).

\paragraph{Not selection.} The daughters' gain is not selection into who is observed, into work or into marriage, and it
does not run through the mother. The father's layoff does not predict whether a daughter is observed with a wage at
25--29 or 45--49, and weighting by the inverse of that probability leaves her slope at 0.176$^{**}$. Her employment does
not change, and the within-person wage slope is 0.169$^{**}$ with partnership and full-year work held fixed. The mother
does not work more in the father's layoff waves (0.011, s.e.\ 0.010), and the daughters' gain is the same whether she did
or not (panel C of Table~\ref{tab:layoff_referee}).\footnote{Code: \texttt{105.DaughtersMechanism.py},
\texttt{109.DaughtersRefereeChecks.py}.}

\paragraph{Schooling and partners.} Two facts point to where the contrast may come from, without settling it. Schooling
moves apart: a teen layoff goes with 0.89 more years of schooling for daughters than for sons ($p = 0.069$), and holding
each child's schooling and its rising return fixed shrinks the within-person difference from 0.254$^{**}$ to 0.162
($p = 0.17$), about a third.\footnote{Code: \texttt{115.LayoffSchooling.py}.} And
daughters of laid-off fathers partner with men who earn more: their partners' earnings rise by 0.677$^{***}$ at 45--49,
which is why their family income rises (panel D of Table~\ref{tab:layoff_mechanism}). The gain does not reach the
daughters' earnings with zero years kept (a change of 0.188, s.e.\ 0.418, in a Poisson model, against $-0.523^{**}$ for
sons): it is a higher wage when they work. Why daughters of laid-off fathers are paid more and partner with higher
earners, the data here do not explain.
